% Бүлэг 3 — Хэрэгжүүлэлт

\chapter{ХЭРЭГЖҮҮЛЭЛТ}
\label{Chapter4}
\pagecolor{white}

%-------------------------------------------------------------------------------

\section{Үүлэн байршуулалт}

\subsection{Docker контейнержуулалт ба AWS EC2}

Node.js/Express бэкэнд сервер нь Docker контейнер хэлбэрээр AWS EC2 (ap-southeast-1, Сингапур) дээр ажиллана. Docker ашиглахын гол зорилго нь хөгжүүлэлтийн болон продакшн орчин яг ижил байлгах явдал юм. Бэкэнд нь \code{backend/Dockerfile}-аар тодорхойлогдсон дүрс ашиглана. Node.js 18 Alpine суурь дүрс сонгосноор эцсийн дүрсний хэмжээ жижиг байж, татах болон байршуулах хугацааг хурдасгана.

\begin{lstlisting}[language={}, caption={backend/Dockerfile --- Node.js 18 Alpine суурьтай бэкэндийн дүрс}]
FROM node:18-alpine
WORKDIR /app
COPY package*.json ./
RUN npm ci --only=production
COPY . .
EXPOSE 3000
CMD ["node", "server.js"]
\end{lstlisting}

EC2 instance дээр Docker суулгасны дараа ECR-аас дүрс татан контейнер ажиллуулна. \code{--restart unless-stopped} тугаар контейнер дахин эхлүүлэгддэг тул EC2 дахин асах үед бэкэнд автоматаар ажиллана. \code{.env} файл нь GROQ\_API\_KEY, EMAIL\_USER, EMAIL\_PASS болон Firebase тохиргооны нууц мэдээллийг хадгалдаг бөгөөд Git-д оруулдаггүй. Дүрслэлийн архитектурыг доорх диаграммаар харуулав.

\begin{figure}[h]
\centering
\begin{tikzpicture}[
  box/.style={draw, rounded corners=6pt, minimum width=3.5cm,
              minimum height=1.1cm, align=center, font=\small, thick, fill=gray!12},
  arr/.style={-latex', thick},
  font=\small
]

\draw[dashed, rounded corners=8pt, fill=gray!5]
  (-1.5,-0.2) rectangle (3.5,1.8);
\node[font=\tiny\bfseries, above] at (1,1.8) {Клиент (Хэрэглэгчийн төхөөрөмж)};
\node[box] (phone) at (1,0.8) {Flutter App\\Android / iOS};

\draw[dashed, rounded corners=8pt, fill=gray!8]
  (4.5,-0.2) rectangle (10.5,3.8);
\node[font=\tiny\bfseries, above] at (7.5,3.8) {AWS Cloud (ap-southeast-1)};
\node[box] (ec2box)  at (7.5,2.8) {AWS EC2 Instance\\t3.micro};
\node[box] (docker)  at (7.5,1.5) {Docker Container\\Node.js/Express :3000};
\node[box] (ecrbox)  at (7.5,0.2) {AWS ECR\\Docker Image Registry};

\draw[dashed, rounded corners=8pt, fill=gray!5]
  (-1.5,-4.8) rectangle (3.5,-2.2);
\node[font=\tiny\bfseries, above] at (1,-2.2) {Google Firebase Cloud};
\node[box] (auth)      at (1,-3.2) {Firebase Auth};
\node[box] (firestore) at (1,-4.2) {Cloud Firestore};

\draw[dashed, rounded corners=8pt, fill=gray!5]
  (4.5,-4.8) rectangle (10.5,-2.2);
\node[font=\tiny\bfseries, above] at (7.5,-2.2) {DevOps / CI-CD Pipeline};
\node[box] (github)   at (7.5,-3.2) {GitHub Repository};
\node[box] (actions)  at (7.5,-4.2) {GitHub Actions Workflow};

\draw[arr] (phone.east) -- node[above,font=\tiny]{REST API} (docker.west |- phone.east);
\draw[arr] (phone.south)-- node[left, font=\tiny]{SDK}      (auth.north);
\draw[arr] (docker.south)--node[right,font=\tiny]{pull image}(ecrbox.north);
\draw[arr] (actions.north)--node[right,font=\tiny]{docker push}(ecrbox.south);

\end{tikzpicture}
\caption{Үүлэн байршуулалтын архитектур --- AWS EC2, Firebase, CI/CD\\{\footnotesize Эх: Зохиогчийн боловсруулалт}}
\label{fig:deployment}
\end{figure}

\subsection{GitHub Actions CI/CD хоолой}

GitHub Actions-ийн CI/CD хоолой нь хөгжүүлэгч \code{backend/} хавтаст өөрчлөлт push хийх бүрт автоматаар ажилладаг. Нийт байршуулалт 3--4 минут болдог бөгөөд хүний гараар сервер дээр нэвтрэх шаардлага байхгүй болдог. Workflow нь 6 алхамтай: checkout, AWS нэвтрэлт тохируулах, ECR нэвтрэх, Docker build хийж ECR-д push хийх, EC2-д SSH-аар холбогдон шинэ дүрс татаж контейнер солих.

\begin{lstlisting}[language={}, caption={GitHub Actions --- автомат байршуулалтын workflow (.github/workflows/deploy.yml)}]
name: Deploy Backend
on:
  push:
    paths: ['backend/**']
    branches: [main]
jobs:
  deploy:
    runs-on: ubuntu-latest
    steps:
      - uses: actions/checkout@v3
      - name: Configure AWS credentials
        uses: aws-actions/configure-aws-credentials@v4
        with:
          aws-access-key-id:     ${{ secrets.AWS_ACCESS_KEY_ID }}
          aws-secret-access-key: ${{ secrets.AWS_SECRET_ACCESS_KEY }}
          aws-region: ap-southeast-1
      - name: Login and push to ECR
        id: login-ecr
        uses: aws-actions/amazon-ecr-login@v2
      - name: Build and push Docker image
        env:
          ECR_REGISTRY: ${{ steps.login-ecr.outputs.registry }}
        run: |
          docker build -t $ECR_REGISTRY/dz-zaal-backend:latest ./backend
          docker push $ECR_REGISTRY/dz-zaal-backend:latest
      - name: Deploy to EC2
        env:
          EC2_HOST: ${{ secrets.EC2_HOST }}
          EC2_SSH_KEY: ${{ secrets.EC2_SSH_KEY }}
          ECR_REGISTRY: ${{ steps.login-ecr.outputs.registry }}
        run: |
          printf '%s\n' "$EC2_SSH_KEY" | tr -d '\r' > /tmp/deploy_key
          chmod 600 /tmp/deploy_key
          ssh -i /tmp/deploy_key -o StrictHostKeyChecking=no \
            ec2-user@$EC2_HOST "
              docker pull $ECR_REGISTRY/dz-zaal-backend:latest
              docker stop dz-zaal-backend || true
              docker rm   dz-zaal-backend || true
              docker run -d --name dz-zaal-backend \
                --restart unless-stopped \
                --env-file ~/backend/.env \
                -p 3000:3000 \
                $ECR_REGISTRY/dz-zaal-backend:latest"
\end{lstlisting}

CI/CD хоолойн дарааллыг доорх диаграммаар харуулав. Git push-ийн дараа Docker build амжилттай болсон тохиолдолд ECR push болж, EC2-д автоматаар байршуулагддаг.

\begin{figure}[h]
\centering
\begin{tikzpicture}[
  cicdstep/.style={draw, rounded corners=5pt, minimum width=3cm,
               minimum height=0.9cm, align=center, font=\small, thick, fill=gray!10},
  dec/.style={draw, diamond, fill=gray!20, aspect=2.2,
              align=center, font=\scriptsize, inner sep=2pt},
  arr/.style={-latex', thick},
  font=\small
]

\node[cicdstep] (push)    at (0,8)    {git push\\backend/};
\node[cicdstep] (trigger) at (0,6.8)  {GH Actions\\Trigger};
\node[cicdstep] (build)   at (0,5.6)  {Docker Build};
\node[dec]      (ok)      at (0,4.2)  {Build\\амжилттай?};
\node[cicdstep] (ecr)     at (0,3.0)  {ECR Push\\(AWS Registry)};
\node[cicdstep] (ssh)     at (0,1.8)  {SSH to EC2};
\node[cicdstep] (pull)    at (0,0.6)  {docker pull\\шинэ image};
\node[cicdstep] (run)     at (0,-0.6) {docker stop old\\docker run new};
\node[cicdstep, fill=gray!25] (fail) at (4,4.2) {Workflow fail\\мэдэгдэл};

\draw[arr] (push)    -- (trigger);
\draw[arr] (trigger) -- (build);
\draw[arr] (build)   -- (ok);
\draw[arr] (ok)      -- node[left,font=\scriptsize]{Тийм}  (ecr);
\draw[arr] (ok)      -- node[above,font=\scriptsize]{Үгүй} (fail);
\draw[arr] (ecr)     -- (ssh);
\draw[arr] (ssh)     -- (pull);
\draw[arr] (pull)    -- (run);

\end{tikzpicture}
\caption{GitHub Actions CI/CD хоолойн дэлгэрэнгүй дарааллын диаграмм\\{\footnotesize Эх: Зохиогчийн боловсруулалт}}
\label{fig:cicd}
\end{figure}

%-------------------------------------------------------------------------------

\section{Туршилт ба үнэлгээ}

\subsection{Функциональ туршилт}

Туршилт нь 2-р бүлгийн 2.1 хэсэгт тодорхойлсон функциональ шаардлага бүрийг баталгаажуулах зорилготой. Туршилтын орчин: Android 13 (Samsung Galaxy A54), iOS 16 (iPhone 13), AWS EC2 t3.micro (ap-southeast-1), Firebase Cloud Firestore (us-central1), 4G LTE болон Wi-Fi сүлжээ.

Апп-ын үндсэн функц болгоныг бодит хэрэглэгчийн нөхцөлд шалгасан бөгөөд тест бүр нь шаардлагын аль хэсэгтэй холбогдохыг тусгав. Давхардсан захиалгыг шалгахдаа хоёр хэрэглэгч яг нэгэн зэрэг ижил цагт захиалах нөхцөлийг тусгайлан туршиж, Firestore-ийн transaction механизм зөв ажилласныг баталгаажуулсан.

\begin{longtable}{|p{0.4cm}|p{2.5cm}|p{3cm}|p{3.2cm}|p{1.5cm}|}
\caption{Функциональ туршилтын тохиолдлууд} \label{tab:functional_tests} \\
\hline
\textbf{\#} & \textbf{Шаардлага} & \textbf{Туршилтын үйлдэл} & \textbf{Хүлээгдэж буй үр дүн} & \textbf{Дүгнэлт} \\
\hline
\endfirsthead
\hline
\textbf{\#} & \textbf{Шаардлага} & \textbf{Туршилтын үйлдэл} & \textbf{Хүлээгдэж буй үр дүн} & \textbf{Дүгнэлт} \\
\hline
\endhead
1 & Email бүртгэл, баталгаажуулалт & Нэр, email, нууц үг оруулж бүртгүүлэх & Баталгаажуулах код email-д ирэх, \code{emailVerified:true} болох & \textcolor{green!60!black}{\textbf{Тэнцсэн}} \\
\hline
2 & Заалны жагсаалт харах & Нүүр дэлгэц нээх & 5 байгууламж, үнэ, байршил зөв гарах & \textcolor{green!60!black}{\textbf{Тэнцсэн}} \\
\hline
3 & Цагийн хуваарь харах & Заал, огноо сонгох & 16 цагийн үе (08:00--24:00) харагдах & \textcolor{green!60!black}{\textbf{Тэнцсэн}} \\
\hline
4 & Цаг захиалах, DBK код авах & Заал, огноо, цаг сонгож баталгаажуулах & DBK-xxxxx кодтой захиалга үүсэх & \textcolor{green!60!black}{\textbf{Тэнцсэн}} \\
\hline
5 & Давхардал хориглох & Захиалагдсан цагт дахин захиалах & ``Уг цаг авагдсан байна'' алдаа гарах & \textcolor{green!60!black}{\textbf{Тэнцсэн}} \\
\hline
6 & Захиалга цуцлах & Идэвхтэй захиалгыг цуцлах & Статус \code{cancelled} болж шинэчлэгдэх & \textcolor{green!60!black}{\textbf{Тэнцсэн}} \\
\hline
7 & Захиалгын түүх харах & Захиалгууд tab руу орох & Захиалга статусаар шүүгдэж харагдах & \textcolor{green!60!black}{\textbf{Тэнцсэн}} \\
\hline
8 & Локал мэдэгдэл авах & Захиалга хийж хүлээх & Захиалгаас 1 цагийн өмнө мэдэгдэл ирэх & \textcolor{green!60!black}{\textbf{Тэнцсэн}} \\
\hline
9 & AI чатбот & ``Ямар заалнууд байна вэ?'' асуух & Монгол хэлээр заалнуудын жагсаалт ирэх & \textcolor{green!60!black}{\textbf{Тэнцсэн}} \\
\hline
10 & Тогтмол захиалга & fixed\_bookings цагийг шалгах & Тэр цагийн үе хаагдсан байдлаар харагдах & \textcolor{green!60!black}{\textbf{Тэнцсэн}} \\
\hline
\end{longtable}

\noindent{\footnotesize Эх: Зохиогчийн боловсруулалт}

Системийн гурван бүрэлдэхүүн хэсгийн хоорондын нэгтгэлийн туршилтыг дараах хүснэгтэд харуулав.

\begin{table}[h]
\centering
\caption{Нэгтгэлийн туршилтын үр дүн}
\label{tab:integration_tests}
\begin{tabularx}{\textwidth}{|p{0.4cm}|X|p{3cm}|p{2cm}|}
\hline
\textbf{\#} & \textbf{Тайлбар} & \textbf{Холбоос} & \textbf{Үр дүн} \\
\hline
1 & Захиалга үүсгэхэд API дуудаж Firestore-д хадгалагдах & Flutter→EC2→Firestore & \textcolor{green!60!black}{\textbf{Тэнцсэн}} \\
\hline
2 & Email бүртгэлийн дараа баталгаажуулах код илгээгдэх & Flutter→Firebase Auth→EC2→Gmail & \textcolor{green!60!black}{\textbf{Тэнцсэн}} \\
\hline
3 & Backend-д push хийхэд Docker build, ECR push, EC2 deploy автоматаар явагдах & GitHub→ECR→EC2 & \textcolor{green!60!black}{\textbf{Тэнцсэн}} \\
\hline
4 & AI чатботын хүсэлт EC2 дамжин Groq API-д хүрч хариулт ирэх & Flutter→EC2→Groq & \textcolor{green!60!black}{\textbf{Тэнцсэн}} \\
\hline
\end{tabularx}\\[4pt]
{\footnotesize Эх: Зохиогчийн боловсруулалт}
\end{table}

3 дугаар тохиолдолд GitHub Actions-ийн \code{ssh-action} нь CRLF мөрийн төгсгөлтэй SSH түлхүүрт алдаа гаргаж байсан тул \code{printf} командыг ашиглан засварлав. Хожим CI/CD хоолой тасралтгүй ажиллаж байгааг 5 дараалсан push-ээр баталгаажуулсан.

\subsection{Гүйцэтгэлийн хэмжилт ба хэрэглэхэд хялбар байдал}

Системийн үндсэн үйлдэл бүрийн хариу цагийг 10 удаа давтан хэмжиж дунджийг гаргасан. Туршилтыг 4G LTE сүлжээн дээр хийсэн бөгөөд хэмжилтийн бүх үзүүлэлт 2.1 хэсэгт тавьсан функциональ бус шаардлагын хязгаарыг хангасан.

\begin{table}[h]
\centering
\caption{Системийн хариу цагийн хэмжилт (секунд) болон шаардлагатай харьцуулалт}
\label{tab:performance}
\begin{tabularx}{\textwidth}{|X|p{2cm}|p{2cm}|p{2cm}|p{2cm}|}
\hline
\textbf{Хэмжигдэхүүн} & \textbf{Дундаж} & \textbf{Бага} & \textbf{Их} & \textbf{Шаардлага} \\
\hline
Апп анх ачаалах & 2.3 & 2.1 & 2.8 & < 3.0 с \\
\hline
Firestore хүсэлт & 0.6 & 0.4 & 0.9 & < 2.0 с \\
\hline
EC2 API хариулт & 0.8 & 0.6 & 1.2 & < 2.0 с \\
\hline
Захиалга үүсгэх & 1.1 & 0.9 & 1.5 & < 2.0 с \\
\hline
Groq Llama хариулт & 0.9 & 0.7 & 1.3 & < 5.0 с \\
\hline
Claude Haiku хариулт & 1.4 & 1.1 & 2.0 & < 5.0 с \\
\hline
\end{tabularx}\\[4pt]
{\footnotesize Эх: Зохиогчийн боловсруулалт}
\end{table}

Апп анх ачаалах үед Firebase SDK болон локал мэдэгдлийн инициализаци давхар явагддаг тул 2.3 секунд зарцуулагдсан. Энэ нь шаардлагын 3.0 секундын доор байна. 15 минут тутмын cron ажлыг шалгахын тулд дуусах цаг нь өнгөрсөн захиалга Firestore-д гараар бүртгэн, ажлын дараагийн гүйцэтгэлийг хүлээсэн. 15 минутын дараа захиалгын статус \code{upcoming}-аас \code{completed} болж шинэчлэгдсэн болохыг Firebase Console дээр баталгаажуулсан.

ӨТДС-ийн 5 оюутан болон Даланзадгад хотын 3 иргэн нийт 8 хүнд аппыг танилцуулж 5 даалгавар гүйцэтгүүлсэн: email бүртгүүлэх, заал хайх, захиалга хийх, захиалга харах, AI чатботоос асуух. Үр дүнг доорх хүснэгтэд харуулав.

\begin{table}[h]
\centering
\caption{Хэрэглэхэд хялбар байдлын үнэлгээний үр дүн}
\label{tab:usability}
\begin{tabularx}{\textwidth}{|X|p{2.8cm}|p{2.5cm}|p{2.8cm}|}
\hline
\textbf{Даалгавар} & \textbf{Амжилттай} & \textbf{Дундаж хугацаа} & \textbf{Тэмдэглэл} \\
\hline
Email бүртгэл/нэвтрэх & 8/8 & 55 сек & Баталгаажуулах код 3--8 сек \\
\hline
Заал хайх & 8/8 & 20 сек & Нүүр дэлгэц шууд ачаалагдсан \\
\hline
Захиалга хийх & 7/8 & 1 мин 15 сек & 1 хүн ``Захиалах'' товчийг эхэнд олоогүй \\
\hline
Захиалга харах & 8/8 & 10 сек & Маш хялбар \\
\hline
AI чатбот & 8/8 & 35 сек & Монгол хэлээр асуух хялбар \\
\hline
\end{tabularx}\\[4pt]
{\footnotesize Эх: Зохиогчийн боловсруулалт}
\end{table}

Дизайн, навигаци, захиалгын процесс, мэдэгдэл, AI чатботын ашигтай байдлыг 1--5 оноогоор үнэлүүлсэн дүнд нийт дундаж оноо \textbf{4.3/5.0} гарсан. Захиалгын процессийн оноо 4.1 харьцангуй бага байсан нь ``Захиалах'' товчийг олоход бага зэрэг хугацаа зарцуулагдсантай холбоотой. Цаашид товчны харагдацыг тодруулах замаар сайжруулах боломжтой. Даланзадгад хотын иргэдийн бүгд захиалга хийх, цуцлах үйлдлийг хугацаандаа зөв хийсэн нь систем орон нутгийн хэрэглэгчдэд тохирсон болохыг харуулсан.

Аюулгүй байдлын шалгалтын хувьд Firebase Auth-ийн хамгаалалт (хэт олон оролдлого хориглох, 10 минутын код хугацаа), API түлхүүрийн нууцлал (GitHub Secrets + EC2 \code{.env}), Firestore Security Rules (\code{bookings}-д зөвхөн нэвтэрсэн хэрэглэгч), өгөгдлийн тусгаарлалт (хэрэглэгч зөвхөн өөрийн UID-тэй захиалгыг харах) зэрэг 5 шалгуурын 4 нь бүрэн хангасан. Нэг шалгуур (HTTP → HTTPS) хэсэгчлэн хангасан бөгөөд цаашдын хөгжлийн шатанд SSL нэмэх шаардлагатай.

\subsection{Туршилтын нийт дүн}

Нийт туршилтын тохиолдлын хамрах хүрээг доорх хүснэгтэд нэгтгэн харуулав.

\begin{table}[h]
\centering
\caption{Туршилтын нийт үр дүн --- шаардлагатай харьцуулалт}
\label{tab:test_summary}
\begin{tabularx}{\textwidth}{|X|p{2.5cm}|p{2.5cm}|p{2.5cm}|}
\hline
\textbf{Туршилтын төрөл} & \textbf{Нийт} & \textbf{Тэнцсэн} & \textbf{Хувь} \\
\hline
Функциональ (шаардлага бүрт нэг тохиолдол) & 10 & 10 & 100\% \\
\hline
Нэгтгэлийн (бүрэлдэхүүн хоорондын) & 4 & 4 & 100\% \\
\hline
Гүйцэтгэлийн (хариу цагийн хязгаар) & 6 & 6 & 100\% \\
\hline
\textbf{Нийт} & \textbf{20} & \textbf{20} & \textbf{100\%} \\
\hline
\end{tabularx}\\[4pt]
{\footnotesize Эх: Зохиогчийн боловсруулалт}
\end{table}

Функциональ туршилт нь 2.1 хэсэгт тодорхойлсон шаардлага бүрийг тусдаа тохиолдлоор шалгасан бөгөөд бүгд амжилттай давсан. Гүйцэтгэлийн хэмжилтийн бүх үзүүлэлт функциональ бус шаардлагын хязгаарыг хангасан. Хэрэглэхэд хялбар байдлын нийт дундаж оноо 4.3/5.0 нь систем нь орон нутгийн иргэдэд тохиромжтой болохоо харуулсан. Туршилтын явцад нэг функциональ алдаа байгаагүй хэдий ч хоёр сайжруулах санал гарсан: захиалгын товчны харагдацыг тодруулах, EC2-ийн холболтод SSL нэмэх. Эдгээрийг цаашдын хөгжлийн шатанд оруулах зорилтот болгосон.
